Let k\ge 2 be an integer.
For x\ge 2, define the property (\*)_{k,x}(n):
\[(\*)_{k,x}(n)\quad\text{means: for each } i\in\{1,2,\dots,k\},\; n+i \text{ is divisible by at least one prime } p\le x.\]
Equivalently, each of the k consecutive integers n+1,\dots,n+k has a prime factor \le x.

Let A_{k,x}:=\{n\in\mathbb N: (\*)_{k,x}(n) \text{ holds}\}.
Define S(k) to be the minimal real x such that A_{k,x} has positive (natural) density, i.e.
\[ \limsup_{N\to\infty} \frac{|A_{k,x}\cap[1,N]|}{N} > 0. \]
(For fixed k, the intended meaning of “positive density set of n” in the problem statement is: the asymptotic density exists and is >0, or at least that the upper density is >0; for our purposes we will show that once A_{k,x} is nonempty it already contains an arithmetic progression and hence has positive density.)

Question: Estimate S(k) as k\to\infty. In particular, is S(k)\ge k^{1-o(1)}?
